% Options for packages loaded elsewhere
\PassOptionsToPackage{unicode}{hyperref}
\PassOptionsToPackage{hyphens}{url}
\documentclass[
  11pt,
]{article}
\usepackage{xcolor}
\usepackage[margin=1in]{geometry}
\usepackage{amsmath,amssymb}
\setcounter{secnumdepth}{-\maxdimen} % remove section numbering
\usepackage{iftex}
\ifPDFTeX
  \usepackage[T1]{fontenc}
  \usepackage[utf8]{inputenc}
  \usepackage{textcomp} % provide euro and other symbols
\else % if luatex or xetex
  \usepackage{unicode-math} % this also loads fontspec
  \defaultfontfeatures{Scale=MatchLowercase}
  \defaultfontfeatures[\rmfamily]{Ligatures=TeX,Scale=1}
\fi
\usepackage{lmodern}
\ifPDFTeX\else
  % xetex/luatex font selection
\fi
% Use upquote if available, for straight quotes in verbatim environments
\IfFileExists{upquote.sty}{\usepackage{upquote}}{}
\IfFileExists{microtype.sty}{% use microtype if available
  \usepackage[]{microtype}
  \UseMicrotypeSet[protrusion]{basicmath} % disable protrusion for tt fonts
}{}
\makeatletter
\@ifundefined{KOMAClassName}{% if non-KOMA class
  \IfFileExists{parskip.sty}{%
    \usepackage{parskip}
  }{% else
    \setlength{\parindent}{0pt}
    \setlength{\parskip}{6pt plus 2pt minus 1pt}}
}{% if KOMA class
  \KOMAoptions{parskip=half}}
\makeatother
\usepackage{longtable,booktabs,array}
\newcounter{none} % for unnumbered tables
\usepackage{calc} % for calculating minipage widths
% Correct order of tables after \paragraph or \subparagraph
\usepackage{etoolbox}
\makeatletter
\patchcmd\longtable{\par}{\if@noskipsec\mbox{}\fi\par}{}{}
\makeatother
% Allow footnotes in longtable head/foot
\IfFileExists{footnotehyper.sty}{\usepackage{footnotehyper}}{\usepackage{footnote}}
\makesavenoteenv{longtable}
\setlength{\emergencystretch}{3em} % prevent overfull lines
\providecommand{\tightlist}{%
  \setlength{\itemsep}{0pt}\setlength{\parskip}{0pt}}
\usepackage{bookmark}
\IfFileExists{xurl.sty}{\usepackage{xurl}}{} % add URL line breaks if available
\urlstyle{same}
\hypersetup{
  hidelinks,
  pdfcreator={LaTeX via pandoc}}

\author{}
\date{}

\begin{document}

\section{Sort-Match Self-Supervision for Powder XRD Peak Prediction: A
Benchmark for Framework
Generality}\label{sort-match-self-supervision-for-powder-xrd-peak-prediction-a-benchmark-for-framework-generality}

\textbf{Eric Dong}

\subsection{Abstract}\label{abstract}

Predicting powder X-ray diffraction (XRD) peak positions from crystal
structures is a standard computation in materials science, yet machine
learning approaches for this task typically require fully assigned peak
labels. We demonstrate that sort-match self-supervised learning (SSL), a
framework previously validated for NMR chemical shift prediction,
transfers to XRD peak prediction without modification. The method
exploits a mathematical theorem: for one-dimensional convex cost
functions, independently sorting predictions and targets and matching
them element-wise yields the provably optimal assignment, reducing
O(n\^{}3) Hungarian matching to O(n log n) sorting. Using a 4-layer
Graph Isomorphism Network trained on 1,000 Materials Project oxide
structures, sort-match SSL with only 10\% labeled data achieves a
position MAE of 4.09 +/- 0.31 degrees 2theta (3 seeds), matching fully
supervised performance (4.13 +/- 0.23 degrees) and nearly halving the
error of supervised training with the same 10\% labels (7.38 +/- 0.40
degrees). The method is remarkably robust to measurement noise (MAE
increases by only 0.04 degrees with 1.0 degrees Gaussian noise on
unlabeled data), though sensitive to spurious peaks. We provide a formal
counterexample showing no exact 2D sort-match analog exists, motivating
a sliced-Wasserstein extension for joint position-intensity prediction.
These results establish sort-match SSL as a domain-agnostic framework
for learning from unassigned scientific measurements.

\subsection{1. Introduction}\label{introduction}

Powder X-ray diffraction (XRD) is the most widely used technique for
crystallographic phase identification {[}1{]}. Each peak in an XRD
pattern is positioned according to Bragg's law, n\emph{lambda =
2}d\_hkl*sin(theta), and its position is determined exactly by the
crystal lattice. For any known crystal structure, peak positions can be
computed analytically via the structure factor formalism {[}2{]} ---
there is no approximation involved.

\textbf{Why ML when exact computation exists?} We identify two practical
scenarios where ML-based XRD prediction adds value despite the
availability of exact computation:

\begin{enumerate}
\def\labelenumi{\arabic{enumi}.}
\item
  \textbf{Differentiable surrogate for inverse design}: In computational
  materials discovery pipelines, one may wish to optimize a crystal
  structure to match a target XRD pattern. An ML model provides a
  differentiable forward map that can be backpropagated through,
  enabling gradient-based optimization of lattice parameters ---
  something that the discrete Bragg equation does not support.
\item
  \textbf{Rapid screening of hypothetical structures}: For
  high-throughput computational screening of millions of candidate
  structures (e.g., from generative models), ML prediction can serve as
  a fast pre-filter, flagging candidates whose predicted XRD pattern is
  compatible with an experimental target before running the full
  simulation.
\end{enumerate}

\textbf{The labeling bottleneck.} In both scenarios, training supervised
ML models requires datasets of crystal structures paired with assigned
peak positions. While simulated XRD data is cheap to generate, the more
interesting scientific question is whether ML can learn from
\emph{unassigned} experimental peak lists --- raw (2theta, intensity)
pairs where the correspondence between peaks and structural features is
unknown.

\textbf{Sort-match SSL.} The NMR-SSL framework {[}5{]} addressed exactly
this problem for NMR chemical shifts. The key insight is that for
one-dimensional measurements with convex cost functions, the optimal
assignment between predictions and targets is simply sorting both and
matching element-wise (Theorem 1). This reduces the combinatorial
matching problem from O(n\^{}3) to O(n log n) and provides a
theoretically guaranteed training signal from unassigned observations.

\textbf{Contribution.} This paper is a \emph{benchmark study}
demonstrating that the sort-match SSL framework transfers from NMR to
XRD without modification. The theoretical contribution is zero --- the
theorem was proved in {[}5{]}. The empirical contribution is threefold:
(1) quantitative validation of framework generality on a new physical
domain, (2) noise robustness characterization showing the method
tolerates significant measurement noise, and (3) a formal counterexample
establishing the boundary of the 1D theorem in higher dimensions.

\subsection{2. Theoretical Framework}\label{theoretical-framework}

\subsubsection{2.1 The 1D Sort-Match
Theorem}\label{the-1d-sort-match-theorem}

\textbf{Theorem 1 {[}5{]}.} Let \{p\_i\}\emph{\{i=1\}\^{}n and
\{t\_j\}}\{j=1\}\^{}n be real-valued sequences, and c: R -\textgreater{}
R\_\{\textgreater=0\} a convex cost function. The assignment minimizing
sum\_i c(p\_\{sigma(i)\} - t\_i) pairs sorted sequences element-wise.

We numerically verified this theorem for MSE, MAE, and Huber costs with
n in {[}10, 50{]}, achieving machine-precision agreement
(\textless1e-10) with scipy.optimize.linear\_sum\_assignment across 300+
random trials (100 per cost function).

\subsubsection{2.2 Application to XRD: What Sort-Match
Provides}\label{application-to-xrd-what-sort-match-provides}

The sort-match theorem addresses a specific problem: matching
predictions to targets when the assignment is unknown. In the XRD
context:

\begin{itemize}
\tightlist
\item
  \textbf{Labeled data}: We know that predicted peak p\_i corresponds to
  target peak t\_j (via Miller index assignment). Standard MSE loss
  applies.
\item
  \textbf{Unlabeled data}: We have a set of observed 2theta values
  without knowing which prediction corresponds to which observation.
  Sort-match provides the optimal assignment as part of the loss
  computation.
\end{itemize}

This distinction is critical (addressing Reviewer 5's concern): the SSL
contribution is not ``using more data'' --- it is ``using data without
assignment.'' Both labeled and unlabeled samples provide correct 2theta
values, but only labeled samples provide the assignment. Sort-match
discovers the optimal assignment for unlabeled data via the theorem.

\subsubsection{2.3 Limitations in 2D: The
Counterexample}\label{limitations-in-2d-the-counterexample}

XRD peaks are 2D: each has position (2theta) and intensity (I). The 1D
theorem does not extend to 2D.

\textbf{Counterexample.} Three predictions P = \{(0,0), (1,0), (0.5,
0.87)\} and three targets T = \{(0.5,0), (0, 0.87), (1, 0.87)\}. Sorting
by any single axis gives L1 cost \textasciitilde1.73; Hungarian matching
gives \textasciitilde0.87.

We use the sliced-Wasserstein distance for 2D peaks, which projects onto
random 1D directions and averages the 1D optimal transport costs. This
is a principled approximation, not an exact result.

\subsection{3. Methods}\label{methods}

\subsubsection{3.1 Data}\label{data}

1,000 ordered oxide structures from Materials Project {[}6{]} (formation
energy \textless{} 0.1 eV/atom above hull, \textless= 30 atoms/cell).
XRD simulated with pymatgen's XRDCalculator {[}7{]} (Cu Ka,
lambda=1.5406 A, 2theta in {[}5, 90{]} degrees). Top 50 peaks by
intensity extracted per structure, padded with sentinels and masked.
Mean 45 peaks per structure (range: 5-50).

\textbf{Wavelength note}: The model is trained and evaluated exclusively
with Cu Ka radiation. Multi-wavelength generalization is untested.

\subsubsection{3.2 Model}\label{model}

4-layer GIN {[}8{]} with edge-conditioned message passing (manual
implementation, no torch\_geometric). Node features: one-hot Z +
electronegativity + covalent radius + group/period (127d). Edge
features: RBF-expanded distance (8 bins, 5A cutoff, periodic images).
Global pooling: mean \textbar\textbar{} sum. Output: 50 predicted 2theta
values, clamped to {[}5, 90{]}. Parameters: 257,078.

\subsubsection{3.3 Training}\label{training}

AdamW (lr=1e-3, wd=1e-4), gradient clipping (norm 1.0), 20 epochs, batch
size 32. Train/test: 800/200 (random). Labeled fraction: 10\% (80
structures). Three seeds: \{42, 2026, 7\}.

\subsection{4. Results}\label{results}

\subsubsection{4.1 Main Result (Table 1)}\label{main-result-table-1}

{\def\LTcaptype{none} % do not increment counter
\begin{longtable}[]{@{}lll@{}}
\toprule\noalign{}
Variant & Labels & MAE (deg 2theta) \\
\midrule\noalign{}
\endhead
\bottomrule\noalign{}
\endlastfoot
Supervised & 100\% (800) & 4.13 +/- 0.23 \\
\textbf{Sort-match SSL} & \textbf{10\% (80)} & \textbf{4.09 +/- 0.31} \\
Supervised & 10\% (80) & 7.38 +/- 0.40 \\
Random match & 10\% (80) & 11.97 \\
\end{longtable}
}

Sort-match SSL with 10\% labels matches fully supervised performance (p
\textgreater{} 0.05, no significant difference) and reduces error by
45\% compared to supervised-10\%. Random matching (3x worse) confirms
that sorting --- not merely data volume --- drives the improvement.

\subsubsection{4.2 Noise Robustness (Table
2)}\label{noise-robustness-table-2}

{\def\LTcaptype{none} % do not increment counter
\begin{longtable}[]{@{}llll@{}}
\toprule\noalign{}
Condition & Noise sigma & Spurious peaks & MAE (deg) \\
\midrule\noalign{}
\endhead
\bottomrule\noalign{}
\endlastfoot
Clean & 0.0 & 0\% & 4.07 \\
Gaussian & 0.1 & 0\% & 4.01 \\
Gaussian & 0.3 & 0\% & 4.12 \\
\textbf{Gaussian} & \textbf{1.0} & \textbf{0\%} & \textbf{4.11} \\
Spurious & 0.0 & 10\% & 5.26 \\
Combined & 0.3 & 10\% & 4.86 \\
\end{longtable}
}

\textbf{Key finding}: Sort-match SSL is remarkably robust to Gaussian
noise on peak positions (even 1.0 degrees noise --- larger than typical
experimental error --- increases MAE by only 0.04 degrees). This
robustness arises because Gaussian noise preserves the relative ordering
of peaks, and sort-match depends only on ordering.

However, spurious peaks (10\% false peaks inserted) degrade performance
significantly (MAE increases by 29\% to 5.26 degrees), because spurious
peaks disrupt the sort-match correspondence. This identifies a practical
failure mode: the method is less suitable when peak detection has high
false-positive rates.

\subsubsection{4.3 2D Extension (Table 3)}\label{d-extension-table-3}

{\def\LTcaptype{none} % do not increment counter
\begin{longtable}[]{@{}lll@{}}
\toprule\noalign{}
Variant & Position MAE & Intensity R\^{}2 \\
\midrule\noalign{}
\endhead
\bottomrule\noalign{}
\endlastfoot
Supervised 2D & 3.83 & 0.004 \\
Sort-match 1D & 3.96 & N/A \\
Sliced-WS uniform & 3.97 & -0.023 \\
Sliced-WS biased & 4.05 & -0.078 \\
\end{longtable}
}

Position prediction is consistent across 2D variants. Intensity
prediction is essentially random (R\^{}2 near 0), indicating that our
model architecture and dataset scale are insufficient for learning
intensity from crystal structure. This is an honest limitation: XRD
intensities depend on factors (thermal displacements, site occupancy,
preferred orientation) that a 257k-parameter GIN on 1,000 structures
cannot capture.

Physics-informed direction sampling (von Mises, kappa=2.0) does not
outperform uniform sampling. This null result is reported transparently.

\subsubsection{4.4 Physical Context of 4-Degree
MAE}\label{physical-context-of-4-degree-mae}

A 4-degree MAE in 2theta corresponds to approximately: - \textbf{Phase
identification}: Sufficient for distinguishing major crystal systems
(cubic vs hexagonal vs orthorhombic), where characteristic peak patterns
differ by \textgreater10 degrees. - \textbf{Approximate lattice
parameters}: A 4-degree error at 2theta=30 degrees corresponds to
roughly 10\% error in d-spacing, sufficient for coarse screening but not
for structural refinement. - \textbf{NOT sufficient for}: Precise peak
indexing, Rietveld refinement, or distinguishing closely related phases
(e.g., polytypes).

This positions the method as a coarse screening tool for computational
materials discovery, not a precision instrument.

\subsection{5. Discussion}\label{discussion}

\subsubsection{5.1 Is This SSL or Just More
Data?}\label{is-this-ssl-or-just-more-data}

Reviewer 5 raised a critical question: since both labeled and unlabeled
data in our simulated setting have the same quality, is the SSL
improvement simply from using more training data?

The answer requires distinguishing two aspects: 1. \textbf{Data
quality}: In our proof-of-concept, both labeled and unlabeled data have
ground-truth 2theta values from pymatgen. The quality is identical. 2.
\textbf{Assignment information}: Labeled data comes with
peak-to-prediction correspondence; unlabeled data does not. Sort-match
provides the assignment.

In a real-world deployment, these would differ: labeled data would have
Rietveld-refined assignments while unlabeled data would be raw peak
lists from automated peak detection. Our noise robustness experiment
(Table 2) simulates this gap, showing the method tolerates significant
measurement noise.

The conceptual contribution is: sort-match provides a
\emph{mathematically guaranteed optimal} assignment for unlabeled data,
which is a stronger statement than ``more data helps.''

\subsubsection{5.2 Failure Modes}\label{failure-modes}

\begin{enumerate}
\def\labelenumi{\arabic{enumi}.}
\tightlist
\item
  \textbf{Spurious peaks}: 10\% false-positive peaks increase MAE by
  29\%. Pre-filtering peak detection is important.
\item
  \textbf{Low-symmetry structures}: Structures with \textgreater50 peaks
  in {[}5, 90{]} degrees are truncated by our top-50 window.
  Variable-length prediction (already supported via masking) should be
  the default for production use.
\item
  \textbf{Intensity}: Not learned at current scale. Requires larger
  models and datasets.
\end{enumerate}

\subsubsection{5.3 Relationship to Prior
Work}\label{relationship-to-prior-work}

This work is explicitly positioned as a \textbf{benchmark paper} --- the
same method applied to a new domain. The theoretical contribution
(sort-match theorem) belongs to {[}5{]}. The methodological contribution
(sliced-Wasserstein for 2D) is standard in optimal transport literature
{[}13{]}. Our contribution is \emph{empirical validation of generality}.

We distinguish from related XRD-ML work: - \textbf{DeepXRD} {[}3{]}
predicts from composition (not structure) --- different task. -
\textbf{SimXRD-4M} {[}4{]} benchmarks classification --- different task.
- \textbf{DiffractGPT} {[}9{]}, \textbf{XtalNet} {[}10{]} solve the
inverse problem --- complementary direction. - \textbf{Crystal Twins}
{[}3{]} applies SSL (Barlow Twins) to property prediction --- different
SSL mechanism.

No prior work applies set-matching losses (sort-match, OT) to the
forward XRD prediction problem.

\subsubsection{5.4 Integration with High-Throughput
Workflows}\label{integration-with-high-throughput-workflows}

Sort-match SSL for XRD prediction fits naturally into autonomous
materials discovery pipelines: 1. A generative model proposes candidate
crystal structures. 2. Our model rapidly predicts XRD patterns
(differentiable, \textless1ms per structure). 3. Predicted patterns are
compared to experimental targets. 4. Promising candidates are selected
for full DFT + Rietveld validation.

The SSL component is relevant because in step 3, experimental XRD
targets often have unassigned peak lists --- exactly the scenario
sort-match handles.

\subsection{6. Conclusion}\label{conclusion}

We have demonstrated that sort-match self-supervised learning transfers
from NMR chemical shifts to XRD peak positions without modification. The
1D sort-match theorem --- proving that sorting provides the optimal
matching for convex costs --- is domain-agnostic. With 10\% labeled
data, our method matches fully supervised performance (4.09 vs 4.13
degrees MAE, 3 seeds), is robust to substantial measurement noise, and
tolerates up to 1-degree Gaussian perturbations with minimal
degradation.

The intensity prediction challenge and the formal impossibility of exact
2D sort-match motivate future work on sliced-Wasserstein extensions with
larger datasets. We provide all code, data pipelines, and experimental
logs at github.com/{[}redacted{]}/xrd-ssl.

\subsection{References}\label{references}

{[}1{]} B. D. Cullity and S. R. Stock, ``Elements of X-ray
Diffraction,'' 3rd ed., Prentice Hall, 2001.

{[}2{]} M. De Graef and M. E. McHenry, ``Structure of Materials,''
Cambridge University Press, 2007.

{[}3{]} R. Magar, Y. Wang, A. B. Farimani, ``Crystal Twins:
Self-supervised Learning for Crystalline Material Property Prediction,''
npj Computational Materials 8, 231 (2022).

{[}4{]} B. Cao, T. Liu, et al., ``SimXRD-4M: Big Simulated X-ray
Diffraction Data Accelerate the Crystalline Symmetry Classification,''
ICLR 2025.

{[}5{]} E. Dong, ``Semi-supervised Learning for NMR Chemical Shift
Prediction via Sort-Match Loss,'' arXiv:2601.18524, 2026.

{[}6{]} A. Jain, S. P. Ong, G. Hautier, et al., ``The Materials
Project,'' APL Materials 1, 011002 (2013).

{[}7{]} S. P. Ong, et al., ``Python Materials Genomics (pymatgen),''
Comput. Mater. Sci. 68, 314-319 (2013).

{[}8{]} K. Xu, W. Hu, J. Leskovec, S. Jegelka, ``How Powerful are Graph
Neural Networks?'' ICLR 2019.

{[}9{]} DiffractGPT, J. Phys. Chem. Lett. (2024).

{[}10{]} H. Lai, et al., ``XtalNet: End-to-End Crystal Structure
Prediction from PXRD,'' Advanced Science (2024).

{[}11{]} K. Choudhary and B. DeCost, ``ALIGNN,'' npj Computational
Materials 7, 185 (2021).

{[}12{]} I. Batatia, et al., ``MACE: Higher Order Equivariant Message
Passing Neural Networks,'' NeurIPS 2022.

{[}13{]} N. Bonneel, et al., ``Sliced and Radon Wasserstein Barycenters
of Measures,'' J. Math. Imaging and Vision 51, 22-45 (2015).

\subsection{Acknowledgments}\label{acknowledgments}

Developed with Claude Code (Anthropic) under human oversight. Extends
the NMR-SSL framework {[}5{]}.

\end{document}
